\section{VC 维把无限假设类变有限}

上一节末尾的结论是：基数是错的计量单位，正确的问题应该问在数据上——一个假设类在 \(n\) 个点上能产生多少种本质上不同的行为。本节把这句话变成定义：生长函数、打散与 VC 维，然后用它们给出不依赖 \(\abs{\mathcal{H}}\) 的泛化界。它同时回答了``无限假设类什么时候可学''和``需要多少样本''。

\subsection{生长函数：\(n\) 个点上的不同行为数}

考虑二分类 \(h:\mathcal{X}\to\{-1,+1\}\)。给定 \(n\) 个点 \(x_1,\dots,x_n\)，每个假设在这批点上只留下一个标记向量

\[
\bigl(h(x_1),\dots,h(x_n)\bigr)\in\{-1,+1\}^n.
\]

在这批点上行为相同的假设，对这批数据而言就是同一个假设。定义生长函数

\[
\Pi_{\mathcal{H}}(n)
=\max_{x_1,\dots,x_n}
\abs{\set{\bigl(h(x_1),\dots,h(x_n)\bigr)\given h\in\mathcal{H}}},
\]

即 \(\mathcal{H}\) 在 \(n\) 个点上最多能产生的不同标记方式数。显然 \(\Pi_{\mathcal{H}}(n)\le2^n\)。若对某 \(n\) 个点等号成立——这批点的 \(2^n\) 种标记全都能被类中假设实现——称 \(\mathcal{H}\) 打散（shatter）这批点。打散是本节的中心概念：它度量假设类在一批点上``拟合任意标记''的能力。

\subsection{VC 维：还能被打散的最大规模}

\(\mathcal{H}\) 的 VC 维 \(d_{\mathrm{VC}}(\mathcal{H})\) 是能被 \(\mathcal{H}\) 打散的点的最大数目：

\[
d_{\mathrm{VC}}(\mathcal{H})
=\max\set{n\given \Pi_{\mathcal{H}}(n)=2^n}.
\]

注意它只要求``存在''一批这样的点，不要求所有 \(n\) 点集合都被打散——VC 维是表达能力的上限声明，不是对任意配置的保证。

例一：实轴上的区间，\(\mathcal{H}=\set{\ind\set{x\in[a,b]}\given a\le b}\)。任取两点 \(x_1<x_2\)，四种标记都能实现：区间盖住谁就给谁 \(+1\)，取一个两点都不盖住的区间则全给 \(-1\)，故 \(d_{\mathrm{VC}}\ge2\)。任取三点 \(x_1<x_2<x_3\)，标记 \((+,-,+)\) 不可能：含两个端点的区间必含中点。故 \(d_{\mathrm{VC}}=2\)。

例二：\(\R^d\) 中的线性分类器 \(h_{w,b}(x)=\sign(w^Tx+b)\)，\(d_{\mathrm{VC}}=d+1\)。下界的直觉：处于一般位置的 \(d+1\) 个点（如 \(\R^2\) 中不共线的三点）总能用半空间实现任意标记。上界的直觉来自 Radon 定理：任意 \(d+2\) 个点都可分成两组，使两组的凸包相交；把两组标以相反符号，任何半空间都无法分开它们——\(\R^2\) 中就是``一点落在另三点的三角形内''与``四点构成凸四边形、取对角标记''两种情形。完整证明是线性代数，此处只需要结论：线性分类器的 VC 维随输入维数线性增长，``参数个数等于 VC 维''在这里碰巧成立。

例三：一个参数也可以有无限 VC 维。考虑

\[
\mathcal{H}=\set{\,x\mapsto\sign\bigl(\sin(\omega x)\bigr)\given \omega>0\,}.
\]

取等比缩小的点 \(x_i=10^{-i}\)，\(i=1,\dots,n\)。对任意标记 \(y_1,\dots,y_n\in\{-1,+1\}\)，令

\[
\omega=\pi\left(\sum_{i=1}^{n}\frac{1-y_i}{2}\,10^{i}+\frac12\right).
\]

考察 \(\omega x_j/\pi\) 的小数部分：\(i>j\) 的项是整数；\(i=j\) 的项贡献 \((1-y_j)/2\)；\(i<j\) 的项总和不超过 \(\frac12\sum_{k\ge1}10^{-k}=1/18\)；末尾的 \(1/2\) 再贡献一个小的正项。于是 \(y_j=+1\) 时小数部分严格落在 \((0,1)\) 内，\(\sin\) 严格为正；\(y_j=-1\) 时小数部分严格落在 \((1,2)\) 内，\(\sin\) 严格为负。故 \(\sign\bigl(\sin(\omega x_j)\bigr)=y_j\) 对每个 \(j\) 成立。\(n\) 是任意的，所以 \(d_{\mathrm{VC}}=+\infty\)——只有一个参数的类可以打散任意大的点集。

三个例子合起来说明：VC 维度量的是表达能力的组合丰富度，不是参数个数，更不是基数。

\subsection{Sauer 引理：从 \(2^n\) 塌缩到多项式}

生长函数的行为有一个干净的二分化：要么对所有 \(n\) 都有 \(\Pi_{\mathcal{H}}(n)=2^n\)（此时 \(d_{\mathrm{VC}}=+\infty\)）；要么一旦 \(n\) 超过 \(d=d_{\mathrm{VC}}(\mathcal{H})\)，生长立刻减速为多项式。Sauer 引理给出精确形式：

\[
\Pi_{\mathcal{H}}(n)
\le\sum_{i=0}^{d}\binom{n}{i}
\le\left(\frac{en}{d}\right)^{d},
\qquad n\ge d.
\]

取对数得 \(\log\Pi_{\mathcal{H}}(n)\approx d\log(n/d)\)。这正是``无限假设类在有限数据上其实很小''的精确含义——类中函数或许不可数，但它们在 \(n\) 个点上的不同行为只有多项式多种。

\subsection{VC 界：用 \(\log\Pi_{\mathcal{H}}\) 替换 \(\log\abs{\mathcal{H}}\)}

把上一节联合界中的 \(\log\abs{\mathcal{H}}\) 换成生长函数的对数（证明中的对称化步骤带来 \(2n\) 与一些常数），就得到 VC 泛化界：以至少 \(1-\delta\) 的概率，对所有 \(h\in\mathcal{H}\) 同时，

\[
R(h)\le\widehat R(h)
+O\left(
\sqrt{\frac{d_{\mathrm{VC}}\log\dfrac{n}{d_{\mathrm{VC}}}+\log\dfrac1\delta}{n}}
\right),
\]

其中 \(d_{\mathrm{VC}}=d_{\mathrm{VC}}(\mathcal{H})\)，也常把 \(d_{\mathrm{VC}}\log(n/d_{\mathrm{VC}})\) 宽松写成 \(d_{\mathrm{VC}}\log n\)。结构与有限类的界完全一致，只是``类的有效大小的对数''从 \(\log\abs{\mathcal{H}}\) 换成了 \(d_{\mathrm{VC}}\log n\)。

它直接回答``需要多少样本''：沿用上一节的 ERM 论证，超出风险不超过两倍界宽，所以要以至少 \(1-\delta\) 的概率把超出风险压到 \(\varepsilon\) 以内，

\[
n\sim\frac{d_{\mathrm{VC}}\log(1/\varepsilon)+\log(1/\delta)}{\varepsilon^2}
\]

量级即足够。样本量随 VC 维近似线性增长，与 \(\abs{\mathcal{H}}\) 无关：VC 维有限的无限类是可学的；VC 维无限的类（如上面的正弦族）在分布无关的意义下不可学——没有任何有限样本能驯服它。这为``什么样的假设类可学''给出了第一个一般性判据：不看大小，不看参数个数，只看 VC 维。

\subsection{证明在做什么：对称化与 Rademacher 复杂度}

本节不展开证明，但值得记住它的两步骨架，因为第二步本身就是现代工具。第一步对称化：引入一份``幽灵样本''，把真风险换成另一份经验风险，分布问题由此变成两份有限样本的差异问题，从而可以在 \(2n\) 个点上枚举行为——生长函数正是在这里登场的。第二步对标记注入随机符号（Rademacher 变量），把``类能实现多少种标记''归约为``类拟合纯随机噪声的能力''；后者就是 Rademacher 复杂度。它把最坏情形的组合计数换成数据依赖的平均拟合噪声能力，因此更紧，也是下一篇范数界的入口。这里只需记住它的角色：把分布问题归约为随机标记问题。
